\chapter{Teilmengen, Verkleben und Erweiterungen}
\label{b08guide:ch:restrictions}

\section{Vollständige Teilmengenkorrespondenz}

\begin{theorem}[Bild und Urbild als zueinander inverse Bijektionen]
\label{b08guide:thm:subsetcorrespondence}
Sei $f\colon A\bij B$. Dann sind
\[
 f_*\colon\mathcal P(A)\bij\mathcal P(B),
 \qquad f^*\colon\mathcal P(B)\bij\mathcal P(A)
\]
zueinander invers, also $(f_*)^{-1}=f^*=(f^{-1})_*$.
Jeder Zielteilmenge $V\subseteq B$ entspricht genau eine
Teilmenge $U\subseteq A$ mit $f[U]=V$, nämlich $U=f^{-1}[V]$.
Insbesondere induziert dieselbe Zuordnung auch Bijektionen
\[
 \begin{aligned}
 \mathcal P(A)\setminus\{\varnothing\}
   &\bij\mathcal P(B)\setminus\{\varnothing\},\\
 \mathcal P(A)\setminus\{A\}
   &\bij\mathcal P(B)\setminus\{B\}.
 \end{aligned}
\]
Ihre Umkehrungen sind jeweils die entsprechenden Einschränkungen von $f^*$.
\end{theorem}

\begin{theorem}[Passende Einschränkungen auf beiden Seiten]
\label{b08guide:thm:restriction}
Seien $f\colon A\bij B$, $U\subseteq A$ und $V\subseteq B$.
Dann sind äquivalent:
\[
 \begin{aligned}
 f|_U\colon U\bij V
 &\quad\Longleftrightarrow\quad f[U]=V\\
 &\quad\Longleftrightarrow\quad U=f^{-1}[V]\\
 &\quad\Longleftrightarrow\quad
   \forall a\in A\ (a\in U\ \Longleftrightarrow\ f(a)\in V).
 \end{aligned}
\]
Wenn diese Bedingungen gelten, ist
\[
 (f|_U)^{-1}=f^{-1}|_V.
\]
Insbesondere ist $f|_{f^{-1}[V]}\colon f^{-1}[V]\bij V$
für jede Teilmenge $V\subseteq B$ bijektiv.
\end{theorem}

\begin{theorem}[Entfernen zusammengehöriger Teilmengen]
\label{b08guide:thm:complementrestriction}
Seien $f\colon A\bij B$ und $U\subseteq A$. Dann ist
\[
 f|_{A\setminus U}\colon A\setminus U\bij B\setminus f[U]
\]
eine Bijektion mit Umkehrfunktion
$f^{-1}|_{B\setminus f[U]}$.
Für $a\in A$ ergibt sich insbesondere
\[
 f|_{A\setminus\{a\}}\colon
 A\setminus\{a\}\bij B\setminus\{f(a)\}.
\]
\end{theorem}

\section{Verkleben und Erweitern}

\begin{theorem}[Verkleben disjunkter bijektiver Stücke]
\label{b08guide:thm:gluing}
Seien $(A_i)_{i\in I}$ und $(B_i)_{i\in I}$ jeweils Familien
paarweise disjunkter Mengen. Sei eine Familie von Bijektionen
$(f_i\colon A_i\bij B_i)_{i\in I}$ gegeben. Dann gilt
\[
 f=\bigcup_{i\in I}f_i\colon
 \bigcup_{i\in I}A_i\bij\bigcup_{i\in I}B_i,
 \qquad f^{-1}=\bigcup_{i\in I}f_i^{-1}.
\]
Die Vereinigungen der Funktionen sind Vereinigungen ihrer Graphen.
Für $I=\varnothing$ entsteht die leere Bijektion.
\end{theorem}

\begin{definition}[Bijektive Erweiterung]
Seien $A\subseteq A'$, $B\subseteq B'$ und $f\colon A\bij B$.
Eine \emph{bijektive Erweiterung} von $f$ auf diese Mengen ist eine
Bijektion $F\colon A'\bij B'$ mit $F|_A=f$.
\end{definition}

\begin{theorem}[Sämtliche bijektiven Erweiterungen]
\label{b08guide:thm:extensions}
Unter den Voraussetzungen der vorstehenden Definition ist
\[
 h\longmapsto f\cup h
\]
eine Bijektion von $\operatorname{Bij}(A'\setminus A,B'\setminus B)$
auf die Menge aller bijektiven Erweiterungen von $f$ nach $A',B'$.
Ihre Umkehrung lautet $F\mapsto F|_{A'\setminus A}$.
Insbesondere existiert eine solche Erweiterung genau dann, wenn
\[
 \operatorname{Bij}(A'\setminus A,B'\setminus B)\ne\varnothing.
\]
Für jedes zulässige $h$ gilt $(f\cup h)^{-1}=f^{-1}\cup h^{-1}$.
\end{theorem}

\begin{theorem}[Erweiterung um ein frisches Paar]
\label{b08guide:thm:freshpair}
Für $f\colon A\bij B$, $a\notin A$ und $b\notin B$ ist
\[
 F=f\cup\{(a,b)\}\colon A\cup\{a\}\bij B\cup\{b\}
\]
die einzige bijektive Erweiterung auf diese beiden Mengen.
Ihre Umkehrfunktion ist $F^{-1}=f^{-1}\cup\{(b,a)\}$.
\end{theorem}

